\section{Threshold integration, moving onsets, and source convergence}
\label{sec:g-integration}
Let $P(\xi)$ be the scatterer perimeter. The collision section measure is
$\dd s\dd p/(2P)$, the mean roof $\pi A/P$, and the stationary collision
intensity $P/(\pi A)$. These identities follow from the same flight-tube
calculation as \eqref{eq:santalo}: irrational torus lines hit the open
scatterer, rational directions form a null angular set, and the positive
minimal gap excludes finite-time accumulation. In particular, in
first-impact endpoint coordinates the full phase measure is
\begin{equation}\label{eq:g-full-flux}
             \frac{-W_{uv}}{2\pi A(\xi)}\dd u\dd v\dd r.
\end{equation}
Indeed, first variation gives $p=-W_u/(\dd s/\dd u)$; the arclength
factor cancels in $|\dd s\dd p|=|W_{uv}|\dd u\dd v$. The positive
relative twist makes the endpoint map injective at fixed $u$. The
localized stationary path is physical by Lemma~\ref{lem:g-channels}.
Thus \eqref{eq:g-full-flux} is not a substituted transverse ensemble.

\begin{lemma}[Uniform smooth radial integration]\label{lem:g-radial}
There are smooth Morse inverses $\Phi_{j,e}(\xi,w)$ on a common ball,
with all fixed derivatives bounded uniformly in $j$, such that
\[
 W_{j,e}(\xi,\Phi_{j,e}(\xi,w))-jg_e=|w|^2/2,
 \qquad D_w\Phi_{j,e}(\xi,0)=H_{j,e}^{-1/2}.
\]
The phase density \eqref{eq:g-full-flux} in $(w,r)$ is
\begin{equation}\label{eq:g-morse-density}
 \frac{a_{j,e}(\xi,w)}{2\pi A(\xi)\sinh(j\gamma_e)}\dd w\dd r,
 \qquad a_{j,e}(\xi,0)=1,
 \qquad a_{j,e}=1+O(|w|)>0,
\end{equation}
with the same derivative bounds. If $a$ is any such uniformly smooth
amplitude, then
\[
 \frac2\pi\int_{|z|<1}(1-|z|^2)a(\xi,\sqrt{2d}\,z)\dd z
            =a(\xi,0)+dH(\xi,d),
\]
where $H$ is smooth for $d\ge0$ and has uniform fixed-order bounds.
No evenness of $a$ or of the scatterer is required.
\end{lemma}
\begin{proof}
For $E=W-jg$, Taylor's integral formula writes
$E(y)=\frac12 y^tM(y)y$, where
$M(y)=2\int_0^1(1-s)D^2E(sy)\dd s$ is uniformly positive near zero.
The map $y\mapsto M(y)^{1/2}y$ has uniformly invertible differential at
zero and small nonlinear differential on a common ball. Its inverse is
obtained by contraction, then differentiated to each finite order.
This is the smooth version of Lemma~\ref{lem:morse}. Inserting the
Jacobian and \eqref{eq:g-ratio} into \eqref{eq:g-full-flux} gives
\eqref{eq:g-morse-density}.

For the last assertion replace $a(w)$ by $(a(w)+a(-w))/2$ in the
integral. Every odd Taylor term vanishes. Taylor expansion to any
prescribed order, with a sufficiently high differentiated remainder,
then yields an expansion in integer powers of $d$. One can also set
$r=\sqrt d$ first: the integral is a smooth even function of $r$ with
uniform bounds. Successive division of its even Taylor remainder by
$r^2$ shows that it is a smooth function of $r^2$ on the right.
Using two more endpoint derivatives for each division proves the
uniform bounds for every fixed derivative of $H$.
\end{proof}

\begin{proof}[Proof of Theorem~\ref{thm:g-stability}]
In a channel's window $jg_e+d$, each of its $j$ complete flights is at
least $g_e$. Hence the allowed first residual time satisfies
$0<r<d-(W-jg_e)<d<g_*$. The preceding roof cannot truncate this interval,
and the next roof cannot fit after the last hit. The direct
first-impact formula of Proposition~\ref{prop:marked-flux} therefore
applies also to this selected channel at its own threshold.
The coercivity \eqref{eq:g-coercive} implies
$|y_0|^2+|y_j|^2+2\sum_{0<i<j}|y_i|^2\le Cd$ on its sublevel.
It places the entire active set inside the common Morse coordinates,
so no chart-boundary term occurs. Its exact probability is
\[
 \frac1{2\pi A\sinh(j\gamma_e)}
 \int_{|w|^2<2d}(d-|w|^2/2)a_{j,e}(\xi,w)\dd w.
\]
Since the quadratic integral is $\pi d^2$, Lemma~\ref{lem:g-radial}
proves the channel assertion, including its derivative bounds.

For the complete event let $t=jg_*+\varepsilon$ with
$0<\varepsilon<\varepsilon_0$. No $j+2$ hits are possible, since
$t<(j+1)g_*$. Lemma~\ref{lem:g-channels} partitions every contributing
$j$-flight segment into exactly one of the six oriented channels.
Its available excess there is $d_{j,e}=\varepsilon-j(g_e-g_*)\le
\varepsilon$. Channels with nonpositive excess contribute zero.
Summing the exact channel formulas gives \eqref{eq:g-competition}.
The minimum-roof bound proves the negative-side assertion.
\end{proof}

\begin{proof}[Proof of Theorem~\ref{thm:g-marked}]
Write the mark on a stationary bridge as $w_{j,e}+V_{j,e}$, subtracting
the normal value separately at each impact. Chord directions and normal
reflections are smooth maps of neighboring transverse coordinates, with
uniformly positive chord denominators. Lemma~\ref{lem:g-relative}
therefore gives $|V_{j,e}|\le C(|u|+|v|)$ and uniform fixed-order
bounds for all its derivatives. The normal states now depend on $\xi$,
but each summand's deviation vanishes at $(u,v)=0$ for every $\xi$;
parameter differentiation preserves its endpoint localization.
The finite chain rule and $\sum_i(\rho^i+\rho^{j-i})\le C$ prove
the bounds. Multiplying the Morse amplitude by $e^{qV_{j,e}}$ and
applying Lemma~\ref{lem:g-radial} proves \eqref{eq:g-marked}.

For the convergence assertion, exact parity counting gives
\begin{equation}\label{eq:g-parity}
 w_{j,e}=(j+1)\bar f_e+
       \one_{\{j\text{ even}\}}(a_{e,0}-a_{e,1})/2.
\end{equation}
For fixed real $q$ the parity factor is bounded above and below by
positive constants, as is the normalized integral of $a_{j,e}e^{qV}$
on a sufficiently small collar. Also $\csch(j\gamma_e)\asymp
2e^{-j\gamma_e}$, uniformly in $j$. Thus a channel's normalized term
is comparable to $e^{j(q\bar f_e-\gamma_e)}$. Positivity rules out
cancellation: equality in \eqref{eq:g-source-domain} gives a
nondecaying term, and a positive exponent gives a growing term.
This proves necessity as well as sufficiency, including the right limit
at zero. On a compact subset of the strict domain there is a common
margin $\delta>0$. Parameter and source derivatives of the leading
factor cost only a polynomial in $j$; offset derivatives act on the
uniform smooth remainder. Every fixed derivative is bounded by
$C(1+j)^m e^{-\delta j}$ and is summable. The vector-source assertion
is the same argument with scalar products.
\end{proof}

\begin{corollary}[Moving threshold interfaces]\label{cor:g-interfaces}
Within the collar of \eqref{eq:g-competition}, the probability is a
$C^1$ function of the smooth physical variables $(\xi,t)$, with smooth
pieces separated by $t=jg_e(\xi)$. For an orientation $e$, along a smooth
scalar path $s\mapsto(\xi(s),t(s))$ entering $d_{j,e}>0$ transversely,
its contribution has second-derivative jump from the inactive to the active side
\begin{equation}\label{eq:g-interface-jump}
       \frac{(t'(s)-j\,\partial_s g_e)^2}{A(\xi)\sinh(j\gamma_e)}.
\end{equation}
Coincident contributions, including the reversed orientation, add with
their respective activation signs.
The marked formula has the additional factor $e^{qw_{j,e}}$.
One-sided and distributional higher derivatives follow by differentiating
the positive-part expression, with polynomial-in-$j$ coefficient bounds.
\end{corollary}
\begin{proof}
Each summand is $d_+^2[C(\xi)+d_+C(\xi)H(\xi,d_+)]$.
It and its first derivatives match at zero. Only
$2C(\partial_s d)^2$ survives in the second-derivative jump.
All other terms contain a factor $d$. Finite addition proves the
claim at intersections. For higher derivatives the ordinary product
and chain rules on each side, and the derivatives of $d_+^2$ in the
distributional sense, give the stated formulas; no differentiation of
$g_*$ is used.
\end{proof}

For a concrete change of minimizing mechanism, take
$h_s(\theta)=R_*+s\cos(2\theta)$. Differentiating the closest distance
at the circular normal chord gives
$g_e'(0)=-2\cos(2\arg e)$: the endpoint variation is the normal support
variation, since the unperturbed endpoint gradient is zero.
The horizontal pair has derivative $-2$ and the other two pairs derivative
$1$. Reflection symmetry keeps those two lengths equal. Consequently
the horizontal pair minimizes for small positive $s$, while the other
two minimize for small negative $s$. Formula \eqref{eq:g-competition}
continues through $s=0$. A nonminimal pair contributes exactly while
$j(g_e-g_*)<\varepsilon$; the parameter transition scale is therefore
$\varepsilon/j$ at a transverse length crossing. This assertion is about
activation of physical channels, not a new family of symbolic itineraries.

\section{Curvatures invisible to all threshold locations}
\label{sec:g-inverse}
The following construction separates geometric uncertainty from its
statistical measurement. Set $\theta_r=r\pi/3$, $r=0,1,2$, and let
\begin{equation}\label{eq:g-isogap-family}
 h_{R,\alpha,\beta,\zeta}(\theta)
   =R+(1-\cos6\theta)b(\theta),\qquad
 b(\theta)=\alpha+\beta\cos2\theta+\zeta\sin2\theta.
\end{equation}
Here $R$ ranges over a compact interval in $(0,1/2)$, and the three
other parameters are sufficiently small. These are smooth centrally
symmetric strictly convex scatterers, generally noncircular. At each
$\theta_r$ and its opposite, $h=R$, $h'=0$, and
\begin{equation}\label{eq:g-contact-curvatures}
       \kappa_r^{-1}=h(\theta_r)+h''(\theta_r)
                     =R+36b(\theta_r).
\end{equation}
The map from $(\alpha,\beta,\zeta)$ to the three $b(\theta_r)$ is
invertible: its rows are $(1,1,0)$, $(1,-1/2,\sqrt3/2)$ and
$(1,-1/2,-\sqrt3/2)$.

\begin{theorem}[Equal gaps and identifiable curvature amplitudes]
\label{thm:g-inverse}
In \eqref{eq:g-isogap-family} all six shortest gaps are exactly
$g=1-2R$, independently of $(\alpha,\beta,\zeta)$. The entire set of
maximal-count threshold locations is therefore the same. Nevertheless,
write
\[
 C_j=\lim_{d\downarrow0}d^{-2}\Pr(N_{jg+d}=j+1),\qquad
 \gamma_r=\arcosh(1+g\kappa_r).
\]
Together with the observed spacing $g$, the unlabelled sequence
$(C_j)_{j\ge1}$ determines the unordered triple
$(\kappa_0,\kappa_1,\kappa_2)$, including multiplicities, and the free
area $A$, without prior knowledge of $A$. It obeys
\begin{equation}\label{eq:g-mobius}
 C_j=\frac1A\sum_{r=0}^2\csch(j\gamma_r),\qquad
 F_j:=\sum_{\substack{k\ge1\\ k\text{ odd}}}\mu(k)C_{kj}
       =\frac2A\sum_{r=0}^2e^{-j\gamma_r},
\end{equation}
where $\mu(k)$ is the arithmetic M\"obius function. Both series involved
in the inversion are absolutely convergent.

In the subfamily $\beta=\zeta=0$, two coefficients already determine
its shape parameter:
\begin{equation}\label{eq:g-two-amplitude-inverse}
 \alpha=\frac1{36}\left\{
       \frac{g}{C_1/(2C_2)-1}-R\right\}.
\end{equation}
Threshold locations alone cannot determine this parameter.
\end{theorem}
\begin{proof}
For small parameters $h+h''$ remains positive, which proves smooth
strict convexity and \eqref{eq:g-contact-curvatures}. For each nearest
lattice vector $e$ at angle $\theta_r$, the two supporting lines have
separation $1-h(\theta_r)-h(\theta_r+\pi)=g$.
Every displacement from the first body to its translate has projection
at least $g$ in direction $e$, and hence length at least $g$.
Since $h'=0$ at the two normals, the supporting points themselves lie
on the axis and attain that distance. More distant translates have
distance at least $\sqrt3-2\sup h>g$ after a uniform decrease of the
parameter neighborhood. Thus the global shortest length is exactly
$g$ and all six channels attain it. The three curvatures can vary
independently by the displayed invertible matrix.

Central symmetry makes both curvatures of the $r$th pair equal.
Theorem~\ref{thm:g-stability} gives the first formula in
\eqref{eq:g-mobius}, with two orientations per pair. Put
$x_r=e^{-\gamma_r}\in(0,1)$. Expanding
$\csch(j\gamma_r)=2\sum_{m\ge1,\ m\text{ odd}}x_r^{mj}$
and interchanging the absolutely convergent sums gives
\[
 F_j=\frac2A\sum_r\sum_{\substack{n\ge1\\n\text{ odd}}}
        x_r^{nj}\sum_{k\mid n}\mu(k)=\frac2A\sum_r x_r^j.
\]
Absolute convergence follows, for example, by bounding the double sum
by $\sum_{n\ge1}n x_{\max}^{nj}<\infty$; the number of divisors
is at most $n$. The divisor identity leaves only $n=1$.

For completeness we give a finite-rank recovery from the transformed
sequence, so uniqueness is not merely asserted. Let $x_1,\ldots,x_d$
be its distinct nodes, $d\le3$, with weights $w_l=2m_l/A>0$, where
$m_l$ are their multiplicities. Then $F_j=\sum_l w_lx_l^j$.
The Hankel matrices
\[
 H^{(0)}=(F_{p+q+1})_{0\le p,q<d},\qquad
 H^{(1)}=(F_{p+q+2})_{0\le p,q<d}
\]
satisfy $H^{(0)}=V\diag(w_lx_l)V^t$ and
$H^{(1)}=V\diag(w_lx_l^2)V^t$, where $V_{pl}=x_l^p$.
The rank of the full Hankel sequence is $d$, and $H^{(0)}$ is positive
definite. Thus the generalized eigenvalues of this pair are exactly
$x_l$. A Vandermonde system using $F_1,\ldots,F_d$ recovers the $w_l$.
Since $\sum_lm_l=3$, we recover
$A=6/(\sum_lw_l)$ and then $m_l=Aw_l/2$.
Finally $\kappa_l=(\cosh(-\log x_l)-1)/g$ proves the curvature claim.
No label ordering or global determination of the boundary is asserted.

If $b\equiv\alpha$, the three curvatures coincide and
$C_1/C_2=2\cosh\gamma=2(1+g/(R+36\alpha))$.
Solving gives \eqref{eq:g-two-amplitude-inverse}. All threshold locations
are independent of $\alpha$ on a nontrivial open interval, while this
ratio varies strictly with it. This proves the claimed nonredundancy.
\end{proof}

The infinite original sequence enters the M\"obius sums; this is an exact
identification theorem, not a finite-noisy-data reconstruction guarantee.
The special two-amplitude formula needs no such infinite transform.
In a circle, by contrast, $R=(1-g)/2$ and $1/C_j$ already satisfies a
second-order recurrence. Corollary~\ref{cor:multiplier} remains correct;
Theorem~\ref{thm:g-inverse} supplies the additional geometric uncertainty
which is absent from that one-parameter model.

\begin{corollary}[Instability selection at equal metric thresholds]
\label{cor:g-selection}
In \eqref{eq:g-isogap-family}, conditional on the maximal-count event
in a common small collar, the mass of a channel with
$\gamma_r>\min_s\gamma_s$ is at most
$C\exp[-j(\gamma_r-\min_s\gamma_s)]$, uniformly in the offset.
Thus channels with the same exact threshold length can have exponentially
different probabilities. The limiting leading-amplitude weights among
minimizing channels are equal.
\end{corollary}
\begin{proof}
The relative remainders in \eqref{eq:g-competition} are bounded above
and below by positive constants. Divide a channel's coefficient by
one with minimal $\gamma$ and use the exponential expression for
$\csch$. At offset zero the equal minimal coefficients give equal
leading-amplitude weights. This last assertion concerns amplitudes;
it does not identify the fixed-positive-offset limiting corrections.
\end{proof}
